\subsection{What are the necessary and relevant concepts, principles, theoretical and design considerations for understanding the coursework and for supporting the correct results?}
\label{sec:nrcp}
\begin{enumerate}
\item Signal operations modify a signal's independent variable, involving transformations like time shifting x(t-2), time reversal x(-t), and time compression x(2t).
\item  Impulse responses define the output behavior of linear time-invariant systems, such as low-pass filters characterized by h_{1}(t) = e^{-10t}u(t) and bandpass filters utilizing a cosine carrier.
\end{enumerate}

\subsection{How does any new component, not covered in previous coursework, function?}
\label{sec:nfxn}
By:
\begin{enumerate}
\item Utilizing the heaviside(t) function to act as a unit step function u(t), which bounds signals mathematically within specific time intervals.
\item Applying the xcorr function to compute the cross-correlation sequence, which is then scaled by the time differential to correctly evaluate auto-correlation.
\end{enumerate}
	
\subsection{What figures, equations, and/or tables could support your answers in Sec.~\ref{sec:nrcp} and Sec.~\ref{sec:nfxn}?}
\begin{enumerate}
\item Figure 1 supports the signal operations concepts by displaying x(t) alongside its shifted, flipped, and compressed variants across six structured subplots.
\item Figure 4 supports the auto-correlation evaluation by plotting the resulting amplitude R_{xx}(\tau) against the lag $\tau$ in seconds.
\end{enumerate}
	
\subsection{Did you cite more than two publications in your answers in Sec.~\ref{sec:nrcp} and Sec.~\ref{sec:nfxn}?}
Yes.
\subsection{Did you cite any online source in your answers in Sec.~\ref{sec:nrcp} and Sec.~\ref{sec:nfxn}?}
No.